\chapter{Permission before value}\label{ch:permission}

\section{A favourable number cannot create stock}
The potential tells us what a specified state change would mean in the
declared account. It does not make that change possible. A source holding
three units cannot supply an isolated four-unit lossless transfer; a road
that is closed cannot deliver medicine because the delivery would be useful.
Before an EBU quote can describe an executable action, a separate physical
question must be answered: \emph{can this complete transformation occur now?}

For one isolated lossless edge $i\to j$, with nonnegative stock and no other
bound, the answer requires $0\leq q\leq x^{\mathrm{b}}_i$. This simple interval already
separates value from permission. From the state $(18,2)$, a transfer of
seventeen units is possible under that physical rule even though it raises
the potential and has negative EBU. A request for nineteen units can be
entered into the quadratic formula, but it lies outside the event's physical
domain. A mathematical value for an impossible endpoint is a counterfactual,
not a permit.

\section{A group can fail even when each child passes}
Suppose one source holds eight units and two actions each request five.
Either action can pass an isolated source check. Their simultaneous group
cannot pass a frozen-source rule because its aggregate request is ten. For
that rule the necessary condition is
\begin{equation}
 \sum_{a:\,\operatorname{src}(a)=i}q_a\leq x^{\mathrm{b}}_i.
\end{equation}
The same principle applies to shared pump capacity, receiving storage and
any other physical limit that a group uses together. Conservation of the
final total does not settle these limits.

Now imagine a chain in which the middle node begins empty, receives four
units and sends four out during one event. Its final stock is zero, not
negative. Yet a warehouse that must receive and inspect stock before
dispatching it may refuse this event, while a continuously operating pipe
could allow a declared flow-through process. The endpoint is identical;
the physical timing rule is not. We cannot settle the difference by
choosing the more attractive EBU calculation.

\begin{cautionbox}{Three questions, not one}
Physical feasibility asks what can occur. EBU asks how the complete accepted
change alters the declared potential. Institutional authority asks who may
decide and on what terms. A mathematical answer to one question does not
silently answer the other two.
\end{cautionbox}

\section{A finite menu is a study design, not a law of nature}
A research implementation may offer only a finite list of candidate
quantities, even if the represented resource is continuously divisible.
That makes the study reproducible; it does not prove that omitted quantities
are physically impossible. A menu should be generated without looking ahead
to which result will favour a chosen actor policy. Otherwise an outcome can
be made to look like a property of EBU when it is really a property of the
menu that was offered.

An admissible set also does not rank its members. Choosing the largest EBU,
choosing uniformly among identified plans, or prioritizing an urgent
delivery are different decision rules. Each can be studied, but none is
hidden inside the definition of a potential. In particular, an essential
negative-EBU action is not physically prohibited just because it is negative.

\section{What a permission record must say}
For an accepted action, record the predecessor state and its epoch, route or
other physical action identifier, requested and accepted quantities, shared
resources, governing constraints and the resulting complete increment. If a
request was reduced, the rule that reduced it belongs in the record. If
only endpoints were checked, the record must not imply that transient path
constraints were checked as well.

This discipline leaves us with a physically meaningful candidate for
valuation. The next chapter can now calculate its exact EBU without asking
the potential to perform the job of a route, a stock certificate or a human
decision-maker.
